% !PRINTED 418
\weberpage{418}
%<<CONT
und daraus ergibt sich, da $x$ und $x'$ positiv sein sollen:
\[
x' = x, \quad y' = y, \quad (\eta = + 2).
\]
Ist aber $-D = 4$, so hat (12) außerdem noch die Lösung
\[
\xi = \pm 1, \quad \eta = 0,
\]
und es folgt aus (13):
\[
\begin{aligned}
y' &= \mp 2\,x, \\
x' &= \pm \frac{y}{2},
\end{aligned}
\]
und das Zeichen von $\eta$ ist so zu bestimmen, daß $x'$ positiv wird. % [sic] Der Druck nennt \eta; da \eta=0 ist und x'=\pm y/2 gilt, bestimmt tatsächlich das Vorzeichen von \xi das Vorzeichen von x'.
Ist endlich $-D = 3$, so haben wir außer (14) die Lösungen
\[
\xi = \pm 1, \quad \eta = \pm 1,
\]
worin zunächst beide Zeichen beliebig sind.
Dann folgt aber aus (13):
\[
\begin{aligned}
\pm 2\,x' &= x \pm y, \\
\pm 2\,y' &= y \mp 3\,x,
\end{aligned}
\]
und man kann die Vorzeichen auf zwei Arten so wählen, daß $x'$ positiv wird. Wir haben also den Satz:

I. \gesperrt{Bei gegebener Gleichung (3) führt jede eigentliche Lösung der Gleichung (8) zu einer Transformation $\left(\begin{smallmatrix}a, & b\\ c, & \partial\end{smallmatrix}\right)$ von der $n$ten Ordnung. Verschiedene Lösungen von (8) führen im allgemeinen zu nicht äquivalenten Systemen $\left(\begin{smallmatrix}a, & b\\ c, & \partial\end{smallmatrix}\right)$. In den beiden Ausnahmefällen $D = -4$ und $D = -3$ führen je zwei oder je drei verschiedene Lösungen von (8) zu äquivalenten Systemen.}

Bezeichnen wir mit $k$ eine Zahl, die im allgemeinen gleich der Zahl der eigentlichen Lösungen von (8) ist, für $D = -4$ und $D = -3$ aber die Hälfte oder ein Drittel der Zahl dieser Lösungen, so können wir das Theorem I. so fassen:

II. \gesperrt{Aus einer Gleichung (3) können wir $k$ nicht äquivalente Systeme $\left(\begin{smallmatrix}a, & b\\ c, & \partial\end{smallmatrix}\right)$ ableiten, die der Gleichung (1) genügen.}

\begin{center}
\textbf{\S\,115. Die singulären Werte der Invariante $j(\omega)$.}
\end{center}

Die Frage, die zunächst unser Interesse in Anspruch nimmt, ist die nach den Werten der Modulfunktionen von $\omega$ für die
%CONT>>
